\section{总结与延伸}

本讲从 2021 年的三篇工作出发，给出了一个仍然适用的音频建模框架：先决定时间分辨率与 representation，再决定预测目标和 context hierarchy，最后用与任务匹配的 metric 评价，并检查模型是否学到合理中间结构。Transformer 只是其中的 sequence operator；真正决定系统性质的，是它与 waveform、spectrogram、codebook、pooling、wavelet 和 decoder 如何连接。

\lecturefigure{slide-47-final-thoughts.jpg}{课堂结论：Transformer 是重大进展，但不应被视为架构演化终点。}{Stanford Online 官方视频 00:47:19--00:48:13。}

\teachervoice{讲者最后说 Transformer 似乎“现在什么都能解”，但希望这不是终点，并邀请大家寻找影响更大的下一种架构。四年后的论文修订与后续音频模型演化，更说明这是一种研究态度，而不是胜利宣言。}

\subsection{把三篇论文压成一个工程循环}

三篇工作可以统一为五步：先识别任务真正需要的时间尺度；选择 waveform、spectrum 或 code representation；为 local 与 global context 分配不同带宽；选择能检验假设而非制造漂亮数字的 metric；最后检查内部表示是否形成可解释结构。这个循环比“使用 Transformer”更可复用。

\begin{table}[H]
\centering
\small
\caption{本讲三条路线的统一比较。}
\begin{tabularx}{\textwidth}{L{0.17\textwidth}L{0.30\textwidth}L{0.18\textwidth}>{\raggedright\arraybackslash}X}
\toprule
路线 & 输入与长程机制 & 主要指标 & 最重要的限制 \\
\midrule
原始波形生成 & 8-bit 离散样本；局部 attention 加压缩 condition & Top-5 样本预测 & 非听感指标；生成误差会累积。 \\
表征学习 & VQ codes 或增强视图；预测/对比目标 & 线性探测准确率 & Tokenizer 与 augmentation 决定可学信息。 \\
音频理解 & 可学习波形 patches；pooling/wavelet 层级 & 多标签 mAP & 数据与基线范围有限；未覆盖流式成本。 \\
\bottomrule
\end{tabularx}
\end{table}

\subsection{常见误区清单}

\begin{warningbox}{不要犯的六个推理错误}
\begin{enumerate}
  \item 把 raw waveform、spectrogram 与 VQ token 都叫“audio tokens”，却不说明 token rate 与信息损失；
  \item 用 top-5 sample accuracy 推断音乐或语音听感；
  \item 用同一个 $O(n^2)$ 标签忽略 local window、downsampling 与 hierarchy 的实际结构；
  \item 把 supervised/unsupervised linear-probe gap 当成所有下游任务的固定差距；
  \item 因 learned filter 像 sinusoid 就宣称模型完整复现了人类听觉；
  \item 把 2021 年 “Adieu” 标题当成对后续所有卷积或混合架构的否定。
\end{enumerate}
\end{warningbox}

\subsection{面向后续课程与现代系统的连接}

这套框架自然延伸到 neural audio codecs、speech units、audio-language pretraining、diffusion/flow generation、streaming recognition 与高效长序列模型。读这些系统时仍应问同一组问题：codec bitrate 是多少，tokenizer 是否丢 phase/prosody，context 在哪个尺度交互，生成器是否并行，评价是 reconstruction、perception、semantic accuracy 还是 human preference。新模型名字不会消除旧约束。

\begin{importantbox}{最终教学结论}
Audio Transformer 的核心不是“attention 也能处理声音”，而是：表示把连续物理信号变成可学习序列；多尺度结构把高频细节与长程语义分层；目标函数决定表示保留什么；指标边界决定我们能声称什么。
\end{importantbox}

\subsection{拓展阅读}

\begin{itemize}
  \item \href{https://arxiv.org/abs/2106.16036}{Verma and Chafe, \emph{A Generative Model for Raw Audio Using Transformer Architectures}}：sample-level Transformer 与 context conditioning。
  \item \href{https://arxiv.org/abs/2010.11459}{Verma and Smith, \emph{A Framework for Generative and Contrastive Learning of Audio Representations}}：contrastive 与 next-code representation learning。
  \item \href{https://arxiv.org/abs/2105.00335v1}{Verma and Berger, \emph{Audio Transformers} v1}：课堂对应的 raw-waveform understanding、pooling、wavelet 与 learned front end。
  \item \href{https://arxiv.org/abs/1609.03499}{van den Oord et al., \emph{WaveNet}}：dilated causal convolution baseline。
  \item \href{https://arxiv.org/abs/1711.00937}{van den Oord et al., \emph{Neural Discrete Representation Learning}}：VQ-VAE 与 codebook learning。
  \item \href{https://arxiv.org/abs/2005.00341}{Dhariwal et al., \emph{Jukebox}}：hierarchical discrete music tokens。
  \item \href{https://arxiv.org/abs/2010.00475}{Fonseca et al., \emph{FSD50K}}：课堂分类实验的数据集来源。
\end{itemize}

\end{document}
